% Transcription — fonds Alexandre Grothendieck, shelfmark 123, batch 3 (pages 41-47).
% Established from the facsimile archives/batches/123/batch-03.pdf (a mirror of
% https://grothendieck.umontpellier.fr/123.pdf), per the skill
% transcribe-grothendieck. The batch file carries no cover sheet: PDF page k is
% folder page 40+k, and the archivists' pencil numbers bottom-left agree
% (« 41 », « 42 », « 47 » legible). The folder's last seven pages.
%
% Pass: Opus 5.5 (claude-opus-5-5), 2026-09-24 — first pass, unchecked against the pages by a human.
%
% All seven pages are his hand; this batch holds none of the typescripts or
% the 1969 letter the inventory lists, so nothing here is summarised as
% someone else's text and nothing is skipped as unrelated paper.
%   - page 41 continues a sentence begun on page 40 (batch 2): cyclic lines
%     and cyclic elements of g, the element sum e_{alpha_i} + e_{-psi}, the
%     normaliser of a cyclic line, D/mu_h = E^1, then « e) » on pairs (T, gamma)
%     with gamma a Coxeter element of N(T). Fast, heavily reworked, several
%     struck passages; the connective prose is often lost.
%   - page 42 is a cover: an otherwise blank sheet inscribed top right
%     « Résumé de quelques résultats de Kostant ». Per hand.md § 5 it gets no
%     \page{}; its words become the section heading of what follows.
%   - pages 43-47 are that résumé, written fair and numbered (i), (i bis),
%     (ii), (iii), (iv): principal homomorphisms G_m -> G and Killing pairs;
%     Kostant homomorphisms mu_h -> G, the Z/hZ-grading of g, the stabiliser
%     W_K; Kostant's theorem on sections of g^1 (semisimple regular / nilpotent),
%     the torsor of cyclic lines under T/K(mu_h); cyclic lines and their
%     normalisers. Page 47 is written on its top third only.
%
% Notation fixed for the batch: \mathfrak{g} for Lie(G) (his « Ɣ »-like
% gothic g), \mathbf{G}_m, \mu_h, \mathfrak{g}^j for the graded pieces, \psi
% for the highest root, \Sigma for the simple roots, \check{V}(\cdot) for the
% projective space of lines (his V with a háček). Interlinear additions of his
% are given in the line with \add{} only where a word is plainly implied;
% otherwise they are set in the text with a \note{} saying they are
% interlinear.
\documentclass[11pt,a4paper]{article}
% Le préambule commun à toutes les transcriptions du fonds Grothendieck.
%
% Il tient en une idée : **l'incertitude se compose, elle ne se résout pas.**
% Une transcription qui lisse un mot douteux a détruit la seule chose pour
% laquelle on la faisait — un lecteur doit pouvoir savoir, à chaque endroit,
% ce qui était sur le feuillet et ce qui a été ajouté. Les six macros qui
% suivent sont donc l'appareil critique, et non de la décoration.
%
% Les mêmes commandes sont reconnues par `scripts/render.mjs`, qui produit la
% vue de lecture du site. Une macro ajoutée ici doit l'être aussi là-bas,
% faute de quoi le rendu échouera bruyamment — ce qui est voulu : un rendu
% silencieusement faux est pire qu'un rendu absent.

\ProvidesPackage{grothendieck}[2026/08/08 Transcriptions du fonds Grothendieck]

\usepackage{amsmath,amssymb,amsthm}
\usepackage{mathrsfs}
\usepackage{stmaryrd}
% Les diagrammes commutatifs sont partout dans ce fonds : c'est la langue
% même de la géométrie algébrique de Grothendieck, et une transcription qui
% les rendrait en prose ne transcrirait rien.
\usepackage{tikz-cd}
% Français par défaut — c'est la langue des feuillets — mais l'anglais est
% chargé aussi : la traduction et le résumé sont des documents anglais, et la
% typographie française (espaces avant les deux-points) y serait une faute.
% Ces documents basculent par \selectlanguage{english} en tête de corps.
\usepackage[english,french]{babel}
\usepackage{fontspec}
\usepackage{geometry}
\geometry{a4paper,margin=25mm,marginparwidth=28mm}
\usepackage{marginnote}
\usepackage{xcolor}
% ulem, not soul. `soul`'s \st and \ul are built on a character-by-character
% scan that XeLaTeX's font machinery breaks: the document fails with an
% undefined control sequence rather than degrading. ulem does the same two jobs
% and is Unicode-engine safe. `normalem` stops it hijacking \emph, which the
% transcriptions use for Grothendieck's own emphasis.
\usepackage[normalem]{ulem}
\usepackage{hyperref}
\hypersetup{colorlinks=true,linkcolor=black,urlcolor=black,pdfborder={0 0 0}}

% --- Métadonnées du lot -----------------------------------------------------
% Reprises telles quelles par le script de rendu, qui les affiche en tête de
% la vue de lecture. Elles fixent ce que le fichier prétend couvrir : sans
% elles, une transcription flotte sans point d'attache dans un fonds de seize
% mille feuillets.

\newcommand{\folder}[1]{\gdef\grothFolder{#1}}
\newcommand{\batch}[1]{\gdef\grothBatch{#1}}
\newcommand{\leaves}[2]{\gdef\grothFirst{#1}\gdef\grothLast{#2}}
\newcommand{\foldertitle}[1]{\gdef\grothFoldertitle{#1}}
\gdef\grothFolder{?}\gdef\grothBatch{?}\gdef\grothFirst{?}\gdef\grothLast{?}\gdef\grothFoldertitle{}

% --- Le résumé --------------------------------------------------------------
%
% Il ouvre la lecture modernisée, sous le titre « Résumé », et il est composé
% en petit. Ce n'est pas un ornement : c'est le seul passage du document qui
% ne soit pas des mathématiques, et un lecteur doit pouvoir le reconnaître
% avant de l'avoir lu — puis le sauter s'il connaît déjà le sujet, ou s'y
% arrêter s'il ne le connaît pas. Le filet qui le suit marque la reprise du
% corps.

\newenvironment{resume}
  {\small\setlength{\parskip}{0.45em}}
  {\par\medskip\hrule\medskip}

% --- Mots-clés ---------------------------------------------------------------
%
% En anglais, en clôture du résumé de la lecture modernisée : le vocabulaire
% moderne sous lequel chercher ce que les feuillets construisent. Le manifeste
% du site (`npm run manifest`) les extrait pour étiqueter la cote entière —
% cette ligne est la source unique des tags, il n'y a pas de fichier à côté.
\newcommand{\keywords}[1]{\par\smallskip\noindent
  {\footnotesize\textsc{Keywords} --- \textit{#1}}\par}

% --- L'appareil critique ----------------------------------------------------

% Le numéro de feuillet, en marge. C'est l'ancre qui permet de régler un
% désaccord sur une formule en regardant *un* feuillet plutôt qu'une chemise
% de six cents. Le site s'en sert aussi pour tourner les pages du fac-similé
% quand on fait défiler la transcription.
\newcommand{\leaf}[1]{%
  \marginnote{\footnotesize\color{gray!70}\textsf{#1}}%
  \label{leaf:#1}\ignorespaces}

% Illisible : on ne devine pas. Un mot inventé qui se lit comme les autres est
% le pire résultat possible.
\newcommand{\ill}{\textcolor{red!60!black}{[\,\ldots\,]}}

% Lecture douteuse : le mot est donné, et signalé comme incertain.
\newcommand{\uncertain}[1]{\uline{#1}}

% Ajout de l'éditeur — une lettre manquante, un mot sous-entendu.
\newcommand{\add}[1]{\textcolor{blue!50!black}{[#1]}}

% Biffé par Grothendieck lui-même. Ce qu'il a rayé fait partie du document :
% une rature dit souvent où la pensée a changé de direction.
\newcommand{\struck}[1]{\sout{#1}}

% Note du transcripteur, sur l'état matériel du feuillet.
\newcommand{\note}[1]{\par\smallskip{\small\color{gray!60!black}\textit{[#1]}}\par\smallskip}

% Note marginale de Grothendieck, distincte du corps du texte.
\newcommand{\marginal}[1]{\par\smallskip\noindent\hspace{1em}%
  {\small\color{gray!60!black}#1}\par\smallskip}

% --- Titre ------------------------------------------------------------------

\AtBeginDocument{%
  \begin{center}
    {\large\bfseries Fonds Alexandre Grothendieck --- cote n\textsuperscript{o}~\grothFolder}\\[2pt]
    {\itshape\grothFoldertitle}\\[4pt]
    {\small feuillets \grothFirst--\grothLast\ (lot \grothBatch)}\\[2pt]
    {\footnotesize\color{gray!60!black}Transcription établie d'après le fac-similé numérisé
     par l'Université de Montpellier.\\
     Les lectures incertaines sont soulignées, les illisibles notées [\,\ldots\,],
     les ajouts entre crochets.}
  \end{center}
  \vspace{1em}}

\endinput


\folder{123}
\batch{3}
\pages{41}{47}
\dating{1969}
\watermark{Édition de démonstration}
\foldertitle{[Autour de Kostant] : tapuscrits (s.d.), notes manuscrites (s.d.), lettre (1969).}

\begin{document}

\page{41}

\note{la page reprend au milieu d'une phrase commencée à la page 40. Écriture
rapide, très reprise ; les mots de liaison se perdent souvent, les formules
passent.}

\ill{} dans $\mathfrak{g}$ sous $G$), \struck{donc les droites}
« plus générales », si $g$ est un \ill{} élément \struck{régulier principal}
de $G$, alors les droites propres sous $g$ de $\mathfrak{g}$ de valeur propre
une racine primitive $h$-ième \uncertain{exacte}, \struck{\ill{}} $g$ est
régulier principal, et les droites en question sont conjuguées. On appelle de
telles droites de $\mathfrak{g}$ \emph{cycliques}, les points correspondants
($\neq 0$) de $\mathfrak{g}$ \emph{cycliques}. On en construit, avec un couple
de Killing donné, en prenant
\[
  \sum_{1}^{r} \lambda_{i}\, e_{\alpha_{i}} + \mu\, e_{-\psi},
\]
les $\lambda_{i}$ et $\mu$ tous $\neq 0$ ($\psi$ la plus haute racine).

Le centralisateur d'un élément cyclique $x$ de $\mathfrak{g}$ est un tore
maximal $T$, \struck{le centralisateur est un groupe \ill{} $T'$ entre $T$ et
$N(T)$ (ce qui \ill{} il dépend de la droite $D$ engendrée par $x$)} ; le
normalisateur de $D$ dans $G$ est un groupe $\widetilde{T}$,
$T \subset \widetilde{T} \subset N(T)$, $\widetilde{T}$ correspondant à
l'élément de Coxeter, donc n'est \emph{pas réduit à} $T$. Il \uncertain{faudrait}
\struck{\ill{}} déterminer les s.-groupes correspondants
\uncertain{d'ordre $h$, multiples de $h$, donc \ill{}} $\widetilde{W}$ de $W$
(\struck{\ill{}} $\widetilde{W} \subset \mathbf{C}^{*}$ \ill{} un \emph{groupe
cyclique} ; \struck{\ill{}}) \uncertain{d'ordre} $h$, grâce à la propriété
maximale de l'élément de Coxeter, \struck{\ill{}} il est engendré par
l'élément de Coxeter. \ill{}
\note{les segments « d'ordre $h$, multiples de $h$ … » et « $h$, grâce à la
propriété maximale de l'élément de Coxeter » sont interlinéaires ; leur point
d'insertion exact dans la phrase n'est pas sûr.}

\struck{\ill{}} Si $D$ est une droite cyclique, alors
$\mathbf{D}/\mu_{h} \simeq \mathbf{E}^{1}$ canoniquement [se vérifie sur un
cas quasi-déployé, et dans ce cas on a un élément privilégié cyclique
\uncertain{canonique} $\sum_{i} e_{\alpha_{i}} + e_{-\psi}$]. Donc
\ill{} les éléments cycliques qui donnent $1$ dans
\struck{$\mathbf{E}$} $\mathbf{E}(D) \simeq \mathbf{E}^{1}$, \ill{} voilà les
éléments \emph{cycliques principaux} (déjà définis \uncertain{à priori} au
choix d'un signe \struck{\ill{} près, si $x$ est principal, \ill{} est une
racine primitive $2h$-ième \ill{}}). Les \uncertain{éléments} forment un
torseur sous $G$, de stabilisateur un tore maximal $T'$.
\struck{Dans le groupe \ill{}, \ill{}} \ill{} élément de la forme
$\sum e_{\alpha_{i}} - e_{-\psi}$, \ill{}
\marginal{Mais peut-on choisir un gén.\ can.\ de $e_{-\psi}$ ? i.e.\ la
\uncertain{puissance} \struck{\ill{}} \ill{} d'un groupe \uncertain{épinglé}
(\uncertain{comparer} \ill{}) \uncertain{implique-t-il} l'identité
\struck{\ill{}} \ill{} $P_{-\psi} \simeq P_{\psi}$ ? \ill{}}
\note{la note marginale, à gauche en bas de page, est écrite en oblique et
reliée par un trait au membre $e_{-\psi}$ ; un « ? » entouré, en face de
« stabilisateur », est relié à la même ligne.}

\marginal{N.B. \struck{\ill{}} centralisateur d'un \ill{} contient un
\uncertain{élément} \uncertain{régulier} \ill{} \uncertain{tore maximal}}
\note{note de la marge gauche, en face de la première proposition sur le
centralisateur. Plus haut dans la même marge, une note écrite verticalement est
entièrement biffée et ne se lit pas.}

e) Considérons les couples $(T, \gamma)$ d'un tore maximal $T$ de $G$ et d'un
élément $\gamma \in N(T)$ dont l'image dans $W$ est un élément de Coxeter.
Deux tels couples sont conjugués [\ill{}]
\note{la page s'arrête sur cet énoncé ; la justification entre crochets ne se
lit pas.}

\section*{Résumé de quelques résultats de Kostant}

\note{titre de sa main, en haut à droite de la page 42, qui ne porte rien
d'autre : c'est une chemise, elle ne reçoit donc pas de \texttt{page}. Les
pages 43 à 47 sont écrites au propre et se lisent presque en entier.}

\page{43}

$G$ groupe semi-simple adjoint sur une base $S$, de type constant.

(i) À tout couple de Killing $(T, B)$ dans $G$, on associe un hom.
$\mathbf{G}_{m} \xrightarrow{\ i\ } G$ par descente à partir du cas où le
système de racines sur $T$ est déployé, donc où
$T \simeq \mathbf{G}_{m}^{\Sigma}$ ($\Sigma$ ens.\ des racines simples),
\struck{quand} d'où
$\operatorname{Hom}(\mathbf{G}_{m}, \uncertain{T}) \simeq \mathbf{Z}^{\Sigma}$ :
on prend l'élément dont toutes les coordonnées sont $1$. Le centralisateur
d'un tel hom.\ (i.e.\ d'un tel sous-tore) \struck{est} est $T$, donc
\ill{} \uncertain{que fournit} le couple de Killing.

On appelle \struck{hom} \emph{hom.\ principal} de $\mathbf{G}_{m}$ dans $G$
un hom.\ tel qu'il existe un couple de Killing auquel il soit associé. Ce
couple est alors unique [$i$ est régulier], et par suite la condition
envisagée sur $i$ est locale pour fppf. Il y a corr.\ biunivoque entre hom.\
principaux $\mathbf{G}_{m} \to G$ et couples de Killing.

(i bis) \struck{Le sous-tore normalisateur de $i(\mathbf{G}_{m})$} Si $i$ est
principal, $i^{-1} = [\lambda \mapsto i(\lambda^{-1})]$ \uncertain{est
principal} aussi, \struck{et} il est associé au couple de Killing opposé. Le
normalisateur de $i(\mathbf{G}_{m})$ est le sous-groupe $N_{1}(T)$ de $N(T)$
image inverse du sous-groupe $\simeq (\mathbf{Z}/2\mathbf{Z})_{S}$ de $W$
engendré par l'élément $w_{(T,B)}$ de $W$ qui envoie $B$ sur le Borel opposé.
\marginal{($r \geqslant 1$ !)}
\note{« $N_{1}(T)$ » est interlinéaire ; « est principal » est ajouté à
l'encre brune au-dessus de la ligne.}

Sous-tore de rang $1$ \emph{principal} : qui loc.\ \uncertain{fppf} est
l'image d'un hom.\ principal $\Longleftrightarrow$ qui sur un revêt.\ étale
\struck{galoisien} principal de groupe $\mathbf{Z}/2\mathbf{Z}$ est
\ill{}.

\page{44}

Il y a correspondance biunivoque entre ces sous-tores et les sections de
$\mathrm{Kill}/(\mathbf{Z}/2\mathbf{Z})$ [N.B. « le » $W$ opère sur
$\mathrm{Kill}$, et dans $W$ il y a « le » élément $w$ d'ordre $2$ qui
transforme la \struck{B} chambre de Weyl (\uncertain{standard}) en son
opposée, c'est par ce $w$ qu'on divise] : couples d'un tore maximal et d'un
couple de Borel opposés le contenant (mais ces Borels n'étant pas néc.\
définis séparément sur $S$).

(ii) Soit $i : \mathbf{G}_{m} \to G$ un hom.\ principal. Il induit un hom.\
$\mu_{h} \to G$ \ill{}, $h$ le nb de Coxeter.
\marginal{$G$ simple ?}

Un hom.\ $K$ est dit hom.\ \emph{de Kostant} si loc.\ fppf (ou ét.\ f., ceci
revient au même, comme on voit aussitôt) il est obtenu à partir d'un hom.\
principal. \struck{Le} Alors $K$ est \emph{régulier}, i.e.\ le centralisateur
connexe de $K$ est un tore maximal. Faisant opérer $\mu_{h}$ sur
$\mathfrak{g} = \operatorname{Lie}(G)$, la graduation de type
$\mathbf{Z}/h\mathbf{Z}$ correspondante de $\mathfrak{g}$ donne comme
composante de degré $0$ une \struck{\ill{}} sous-algèbre de Cartan
$\mathfrak{t} = \operatorname{Lie}(T)$, $T = \operatorname{Cent}^{0}(i_{h})$,
et la composante de degré $j$, $1 \leqslant j \leqslant h-1$, \struck{est
égale} a comme rang le nombre de racines positives dont l'ordre est
$\equiv j \pmod{h}$, i.e.\ dont l'ordre est $j$ ou $j - h$ ; en fait, si $K$
est associé à un couple de Killing, \struck{déployé} avec syst.\ de racines
associé à $T$ déployé, alors

\page{45}

la composante $\mathfrak{g}^{j}$ ($1 \leqslant j \leqslant h$) de
$\mathfrak{g}$ relative au caractère $\chi^{j}$ est
\[
  \sum_{\substack{\alpha \in R\\ o(\alpha) = j}} \mathfrak{g}_{\alpha}
  \;+\; \sum_{\substack{\alpha \in R\\ o(\alpha) = j-h}} \mathfrak{g}_{\alpha}.
\]
En particulier $\mathfrak{g}^{1}$ est de rang $r + 1$ ($r$ = rang).

Le centralisateur $G^{K} = N(T)^{K}$ de $K$ est un sous-groupe fermé de
$N(T)$, image inverse du sous-groupe \struck{$W$} $W_{K}$ de $W$,
\struck{centralisateur} stabilisateur de $K$ dans $W$. Kostant dit que le rang
de $W_{K}$ est égal à celui du centre de $\widetilde{G}$, et on peut espérer
qu'il soit isomorphe à $\overline{P/R}$, $P = \widehat{T}$ = réseau des
poids, $R$ = réseau des racines, \struck{\ill{}} ou à son dual
\struck{\ill{}} à valeurs dans $\mathbf{Q}/\mathbf{Z}$.
\note{« $G^{K} = N(T)^{K}$ » est interlinéaire.}

La donnée d'un
\[
  K : \mu_{h} \longrightarrow T
\]
de Kostant dans un tore maximal \uncertain{défini par} le système de racines
$R(T)$ de ce \ill{} équivaut à la donnée d'un système de racines $\Sigma'$
\uncertain{$\mathbf{Q}$-simples}, \ill{} \uncertain{$\mathbf{Q}$-base} ;
l'ens.\ des racines $\alpha$ telles que $\mathfrak{g}_{\alpha} \subset
\mathfrak{g}^{1}$ ($\mathfrak{g}^{1}$ désignant la composante de
$\mathfrak{g}$ relative au car.\ $\chi_{1}$ de $\mu_{h}$). En fait, $K$ est
connu quand on connaît \struck{\ill{}} $\Sigma' \subset R$, \struck{\ill{}}
si $h > 2$ \ill{}.
\note{« $R(T)$ de ce » et « $\mathbf{Q}$-simples » sont interlinéaires,
soulignés d'un trait qui les rattache à la ligne ; la fin du paragraphe porte
un passage encadré et biffé, illisible.}

Je vérifie aussi, dans le cas déployé, que le \struck{centralisateur}
stabilisateur $W_{K}$ de $K$ dans $W$ est égal (a priori il est $\subset$) au
stabilisateur de $\Sigma'$ dans $W$, i.e.\ l'ensemble des $w \in W$ tels que
$w(\Sigma) \subset \Sigma'$, i.e.\ tels que $w$ transforme toutes les racines
de $\Sigma$ en

\page{46}

racines de $\Sigma$, sauf une qui est transformée en $-\psi$ ($\psi$ = plus
grande racine).

Déterminer aussi le normalisateur de $K$, et l'action de
$\operatorname{Norm}(K)/\operatorname{Zent}(K)$ sur $\mu_{h}$ ;
\uncertain{transitive dans}
$\operatorname{Aut}(\mu_{h}) = (\mathbf{Z}/h\mathbf{Z})^{*}$ ?

(iii) Soit \struck{\ill{}} $K : \mu_{h} \to T$ un hom.\ principal. Considérons
comme dessus
$\mathfrak{g}^{1}$ $\bigl\{= \sum_{\alpha \in Q} \mathfrak{g}^{\alpha}$ (qui
dépend de $K$)$\bigr\}$. Kostant prouve (en car.\ $0$) :

a) Toute section $\sum x^{\alpha} \in \Gamma\mathfrak{g}^{1}$
($x^{\alpha} \in \Gamma\mathfrak{g}^{\alpha}$) telle que les $x^{\alpha}$ ne
s'annulent pas est \emph{semi-simple régulière}.
\note{il écrit « ½ simple » pour \emph{semi-simple}.}

b) Si un des $x^{\alpha}$ est nulle, elle est \emph{\uncertain{nilpotente}}.

c) Soient deux droites $D, D' \subset \check{V}(\mathfrak{g}^{1})^{ss}$, qui
soient \uncertain{loc.\ \struck{\ill{}}} dans cas a).
Alors $\exists\, t \in \Gamma T$ telle que $D' = \operatorname{ad}(t).D$
[sans unicité : le stabilisateur dans $T$ d'une $D$ comme dessus est
$K(\mu_{h})$, donc $t$ est une section canonique de $T/K(\mu_{h})$.
\ill{} aussi que dans \struck{\ill{}} \ill{}
\struck{$\operatorname{Isom}_{\underline{0}}(D \xrightarrow{\sim} D')$ \ill{}}
$\mu_{h}$-torseur \ill{} la fibre des $t \in T$ \ill{} une canonique
(savoir \struck{\ill{}} la \ill{} qui envoie $D$ dans $D'$), et pour un
\ill{} donné de $D$ avec $D'$ dans cette classe, il $\exists!$ (unique)
$t \in T$ qui l'induit [cette information résulte du fait que \ill{} de $T$
qui \ill{} l'identité sur $D$ est la section $1$]].
\note{le dernier crochet est ajouté à l'encre brune. Dans la marge gauche,
en face de a) à c), une note oblique au crayon ne se lit pas.}

\page{47}

(iv) \emph{Droites cycliques} $D \subset \check{V}(\mathfrak{g})$ : qui loc.\
fppf sont propres de caractère $\chi_{1}$ sous
$\mu_{h} \xrightarrow{\ K\ } G$, $K$ étant de Kostant. Les droites cycliques
sont donc conjuguées. Si $D$ est une droite cyclique, alors (sous réserve que
ce soit \uncertain{unique}) $\operatorname{Zent}^{0}(D)$ est un tore maximal
$T'$, et $\operatorname{Norm}(D)$ est un sous-groupe de $N' = N(T')$
contenant $T'$, qui est l'image inverse d'un sous-groupe de
$W' = N(T')/T'$ qui est \struck{étale d'ordre $h$ sur $S$, et à fibres
cycliques} \struck{fini} can.\ isom.\ à $\mu_{h}$ (dans le \ill{} sur $S$) et
dont un générateur correspond à un élément de Coxeter.
\note{« fini » (biffé) puis « can.\ isom.\ à » sont interlinéaires, au-dessus
du passage biffé. Le reste de la page est blanc.}

\end{document}
